% Options for packages loaded elsewhere
\PassOptionsToPackage{unicode}{hyperref}
\PassOptionsToPackage{hyphens}{url}
\PassOptionsToPackage{dvipsnames,svgnames,x11names}{xcolor}
\documentclass[
  11pt,
]{article}
\usepackage{xcolor}
\usepackage{amsmath,amssymb}
\setcounter{secnumdepth}{-\maxdimen} % remove section numbering
\usepackage{iftex}
\ifPDFTeX
  \usepackage[T1]{fontenc}
  \usepackage[utf8]{inputenc}
  \usepackage{textcomp} % provide euro and other symbols
\else % if luatex or xetex
  \usepackage{unicode-math} % this also loads fontspec
  \defaultfontfeatures{Scale=MatchLowercase}
  \defaultfontfeatures[\rmfamily]{Ligatures=TeX,Scale=1}
\fi
\usepackage{lmodern}
\ifPDFTeX\else
  % xetex/luatex font selection
\fi
% Use upquote if available, for straight quotes in verbatim environments
\IfFileExists{upquote.sty}{\usepackage{upquote}}{}
\IfFileExists{microtype.sty}{% use microtype if available
  \usepackage[]{microtype}
  \UseMicrotypeSet[protrusion]{basicmath} % disable protrusion for tt fonts
}{}
\makeatletter
\@ifundefined{KOMAClassName}{% if non-KOMA class
  \IfFileExists{parskip.sty}{%
    \usepackage{parskip}
  }{% else
    \setlength{\parindent}{0pt}
    \setlength{\parskip}{6pt plus 2pt minus 1pt}}
}{% if KOMA class
  \KOMAoptions{parskip=half}}
\makeatother
\usepackage{color}
\usepackage{fancyvrb}
\newcommand{\VerbBar}{|}
\newcommand{\VERB}{\Verb[commandchars=\\\{\}]}
\DefineVerbatimEnvironment{Highlighting}{Verbatim}{commandchars=\\\{\}}
% Add ',fontsize=\small' for more characters per line
\usepackage{framed}
\definecolor{shadecolor}{RGB}{255,255,255}
\newenvironment{Shaded}{\begin{snugshade}}{\end{snugshade}}
\newcommand{\AlertTok}[1]{\textcolor[rgb]{0.75,0.01,0.01}{\textbf{\colorbox[rgb]{0.97,0.90,0.90}{#1}}}}
\newcommand{\AnnotationTok}[1]{\textcolor[rgb]{0.79,0.38,0.79}{#1}}
\newcommand{\AttributeTok}[1]{\textcolor[rgb]{0.00,0.34,0.68}{#1}}
\newcommand{\BaseNTok}[1]{\textcolor[rgb]{0.69,0.50,0.00}{#1}}
\newcommand{\BuiltInTok}[1]{\textcolor[rgb]{0.39,0.29,0.61}{\textbf{#1}}}
\newcommand{\CharTok}[1]{\textcolor[rgb]{0.57,0.30,0.62}{#1}}
\newcommand{\CommentTok}[1]{\textcolor[rgb]{0.54,0.53,0.53}{#1}}
\newcommand{\CommentVarTok}[1]{\textcolor[rgb]{0.00,0.58,1.00}{#1}}
\newcommand{\ConstantTok}[1]{\textcolor[rgb]{0.67,0.33,0.00}{#1}}
\newcommand{\ControlFlowTok}[1]{\textcolor[rgb]{0.12,0.11,0.11}{\textbf{#1}}}
\newcommand{\DataTypeTok}[1]{\textcolor[rgb]{0.00,0.34,0.68}{#1}}
\newcommand{\DecValTok}[1]{\textcolor[rgb]{0.69,0.50,0.00}{#1}}
\newcommand{\DocumentationTok}[1]{\textcolor[rgb]{0.38,0.47,0.50}{#1}}
\newcommand{\ErrorTok}[1]{\textcolor[rgb]{0.75,0.01,0.01}{\underline{#1}}}
\newcommand{\ExtensionTok}[1]{\textcolor[rgb]{0.00,0.58,1.00}{\textbf{#1}}}
\newcommand{\FloatTok}[1]{\textcolor[rgb]{0.69,0.50,0.00}{#1}}
\newcommand{\FunctionTok}[1]{\textcolor[rgb]{0.39,0.29,0.61}{#1}}
\newcommand{\ImportTok}[1]{\textcolor[rgb]{1.00,0.33,0.00}{#1}}
\newcommand{\InformationTok}[1]{\textcolor[rgb]{0.69,0.50,0.00}{#1}}
\newcommand{\KeywordTok}[1]{\textcolor[rgb]{0.12,0.11,0.11}{\textbf{#1}}}
\newcommand{\NormalTok}[1]{\textcolor[rgb]{0.12,0.11,0.11}{#1}}
\newcommand{\OperatorTok}[1]{\textcolor[rgb]{0.12,0.11,0.11}{#1}}
\newcommand{\OtherTok}[1]{\textcolor[rgb]{0.00,0.43,0.16}{#1}}
\newcommand{\PreprocessorTok}[1]{\textcolor[rgb]{0.00,0.43,0.16}{#1}}
\newcommand{\RegionMarkerTok}[1]{\textcolor[rgb]{0.00,0.34,0.68}{\colorbox[rgb]{0.88,0.91,0.97}{#1}}}
\newcommand{\SpecialCharTok}[1]{\textcolor[rgb]{0.24,0.68,0.91}{#1}}
\newcommand{\SpecialStringTok}[1]{\textcolor[rgb]{1.00,0.33,0.00}{#1}}
\newcommand{\StringTok}[1]{\textcolor[rgb]{0.75,0.01,0.01}{#1}}
\newcommand{\VariableTok}[1]{\textcolor[rgb]{0.00,0.34,0.68}{#1}}
\newcommand{\VerbatimStringTok}[1]{\textcolor[rgb]{0.75,0.01,0.01}{#1}}
\newcommand{\WarningTok}[1]{\textcolor[rgb]{0.75,0.01,0.01}{#1}}
\usepackage{longtable,booktabs,array}
\usepackage{calc} % for calculating minipage widths
% Correct order of tables after \paragraph or \subparagraph
\usepackage{etoolbox}
\makeatletter
\patchcmd\longtable{\par}{\if@noskipsec\mbox{}\fi\par}{}{}
\makeatother
% Allow footnotes in longtable head/foot
\IfFileExists{footnotehyper.sty}{\usepackage{footnotehyper}}{\usepackage{footnote}}
\makesavenoteenv{longtable}
\setlength{\emergencystretch}{3em} % prevent overfull lines
\providecommand{\tightlist}{%
  \setlength{\itemsep}{0pt}\setlength{\parskip}{0pt}}
% Publication template for GS-002.
%
%  Typesetting is not evidence for any claim.
%
% Built 2026-09-08.
%
% xelatex, not pdflatex: the manuscript uses unicode characters that pdflatex
% cannot set without per-character declarations.

\usepackage{fontspec}
\usepackage{unicode-math}
\setmainfont{Latin Modern Roman}
\setmonofont{Latin Modern Mono}[Scale=0.85]
\setmathfont{Latin Modern Math}

\usepackage{amsmath,amssymb,amsthm}
\usepackage{newunicodechar}
\newcommand{\um}[1]{\ensuremath{#1}}
\newunicodechar{ℝ}{\um{\mathbb{R}}}
\newunicodechar{⁰}{\um{^0}}
\newunicodechar{₀}{\um{_0}}
\newunicodechar{₁}{\um{_1}}
\newunicodechar{₂}{\um{_2}}
\newunicodechar{¹}{\um{^1}}
\newunicodechar{²}{\um{^2}}
\newunicodechar{∘}{\um{\circ}}
\newunicodechar{δ}{\um{\delta}}
\newunicodechar{⟨}{\um{\langle}}
\newunicodechar{⟩}{\um{\rangle}}
\newunicodechar{φ}{\um{\varphi}}
\newunicodechar{ᵀ}{\um{^\top}}
\newunicodechar{𝛷}{\um{\Phi}}
\newunicodechar{Φ}{\um{\Phi}}
\newunicodechar{Γ}{\um{\Gamma}}
\newunicodechar{α}{\um{\alpha}}
\newunicodechar{β}{\um{\beta}}
\newunicodechar{κ}{\um{\kappa}}
\newunicodechar{∂}{\um{\partial}}
\newunicodechar{∈}{\um{\in}}
\newunicodechar{−}{\um{-}}
\newunicodechar{∞}{\um{\infty}}
\newunicodechar{≠}{\um{\neq}}
\newunicodechar{≤}{\um{\le}}
\newunicodechar{≥}{\um{\ge}}
\newunicodechar{⟺}{\um{\iff}}
\newunicodechar{→}{\um{\to}}
\newunicodechar{─}{-}
\newunicodechar{│}{|}
\newunicodechar{┌}{+}
\newunicodechar{┐}{+}
\newunicodechar{└}{+}
\newunicodechar{┘}{+}
\newunicodechar{┬}{+}
\newunicodechar{┴}{+}
\newunicodechar{▼}{v}
\newunicodechar{ℓ}{\um{\ell}}
\newunicodechar{Ω}{\um{\Omega}}
\newunicodechar{Σ}{\um{\Sigma}}
\newunicodechar{·}{\um{\cdot}}
\newunicodechar{×}{\um{\times}}
\newunicodechar{⊂}{\um{\subset}}
\newunicodechar{ₙ}{\um{_n}}
\newunicodechar{∎}{\um{\blacksquare}}
\newunicodechar{∼}{\um{\sim}}

\usepackage[margin=1in]{geometry}
\usepackage{longtable,booktabs,array}
\usepackage{microtype}
\usepackage{enumitem}
\setlist{nosep,leftmargin=1.4em}

\usepackage[dvipsnames]{xcolor}
\usepackage[colorlinks=true,linkcolor=RoyalBlue,urlcolor=RoyalBlue,
            citecolor=RoyalBlue,breaklinks=true]{hyperref}
\usepackage{fvextra}
\usepackage{framed}
\fvset{breaklines=true,breakanywhere=true,fontsize=\small}

\AtBeginEnvironment{longtable}{\small}
\setlength{\tabcolsep}{4pt}
\renewcommand{\arraystretch}{1.15}

\providecommand{\tightlist}{%
  \setlength{\itemsep}{0pt}\setlength{\parskip}{0pt}}

\definecolor{quoterule}{gray}{0.75}
\renewenvironment{quote}
  {\begin{leftbar}\small}
  {\end{leftbar}}

\usepackage[htt]{hyphenat}
\usepackage{xurl}
\emergencystretch=3em
\sloppy

\newenvironment{nonclaim}{\begin{quote}\textbf{Non-claim:\ }}{\end{quote}}
\newenvironment{claimblock}{\begin{quote}}{\end{quote}}
\usepackage{bookmark}
\IfFileExists{xurl.sty}{\usepackage{xurl}}{} % add URL line breaks if available
\urlstyle{same}
\hypersetup{
  pdftitle={SO(k)-Invariant Gauge Rigidity for Declared Euclidean Stalks: Local l2 Selection and Outer l1 Aggregation (GS-002)},
  pdfauthor={Jeremy H. Carroll},
  colorlinks=true,
  linkcolor={RoyalBlue},
  filecolor={Maroon},
  citecolor={Blue},
  urlcolor={Blue},
  pdfcreator={LaTeX via pandoc}}

\date{2026-09-21}

\begin{document}
\begin{center}
{\LARGE \textbf{SO(k)-Invariant Gauge Rigidity for Declared Euclidean Stalks}}\\[0.4em]
{\large \textbf{Local \(\ell^2\) Selection and Outer \(\ell^1\) Aggregation (GS-002)}}\\[1em]
\textbf{Author:} Jeremy H. Carroll \quad \textbf{Date:} 2026-09-21\\[0.4em]
\textbf{Version:} v2.0.0 Pre-Print \quad \textbf{License:} CC-BY-4.0
\end{center}
\vspace{1em}

\subsection{Abstract}\label{abstract}

We study vector-valued defect diagnostics on finite directed graphs with
declared Euclidean stalks. The classical sphere-transitivity argument
shows that a faithful absolutely homogeneous
\(\mathrm{SO}(k)\)-invariant local gauge is a positive multiple of
\(\|\cdot\|_2\). Substitution into a separately declared strong edgewise
sum gives \(J_{2,1,c}(\omega)=\sum_e c_e\|\omega_e\|_2\); symmetry alone
does not determine the outer aggregation. Edge exchange symmetry makes
weights constant on declared orbits, while unit normalization fixes them
to one.

For identity transports, finite-dimensional convex duality identifies
the quotient norm modulo vertex repairs with a maximum over
divergence-free vector witnesses satisfying \(\|\varphi_e\|_2\le c_e\).
We give an attained single-cycle formula
\((\min_e c_e)\|S(\omega)\|_2\), a normalized triangle separation of
\(2\) versus \(\sqrt2\), and a closure bound under \(d_1d_0=0\). These
are applications of standard norm, cycle-space, and convex-duality
results with explicit observer hypotheses. The companion Lean
development mechanizes finite and abstract kernels, including planar
rotation transitivity and weak duality; general \(\mathrm{SO}(k)\)
transitivity and strong duality remain prose or imported steps. Python
examples evaluate analytical formulas and numerical certificates in
floating-point arithmetic. Finite here describes the carrier and stalk
dimension; the Euclidean scalar and symmetry structures are declared
continuum inputs.

\begin{quote}
\textbf{What this does not claim.} This paper does not claim global
Euclidean physics, global \(\ell^1\) geometry, unweighted
\(\ell^1(E;\ell^2)\) without edge normalization, or that every
vector-valued observer must use \(\mathrm{SO}(k)\). It does not claim
that the triangle example shows the general theory by itself. It claims
only that, under declared \(\mathrm{SO}(k)\)-invariant local gauge
symmetry and exact edge-additive replica extensivity, the licensed
finite diagnostic is weighted \(\ell^1(E;\ell^2)\), with unweighted form
requiring normalization or edge-uniformity.

\textbf{Why a researcher should care.} Many finite obstruction systems
have vector-valued local defects: displacement fields, force/strain
residuals, multi-channel sensor mismatches, projected dynamics errors,
and vector-valued presheaf gluing failures. Treating such a defect as
independent scalar coordinates gives a different diagnostic from
treating it as one Euclidean vector. This paper identifies the precise
observer assumption that justifies the Euclidean treatment: continuous
\(\mathrm{SO}(k)\) rotational invariance of the local gauge. It then
separates that local symmetry choice from the outer edge aggregation
assumption, showing that \(\ell^1(E;\ell^2)\) is a conditional observer
geometry, not a global metric.
\end{quote}

\subsection{1. Introduction and Structural
Placement}\label{introduction-and-structural-placement}

The classification of diagnostic geometries over finite graphs forms the
basis of categorical obstruction analysis. Previous classification
theorems show that when an observer's operational domain demands
edge-decomposable aggregation and maximal bounded-linear compatibility
on scalar stalks, the resulting diagnostic is a weighted \(\ell^1\)
coboundary norm, becoming uniform only under edge-uniformity
assumptions.

However, moving from scalar signals to Euclidean geometries requires
stating the structural constraints governing vector variables.
Evaluating a 2D vector defect \(\vec{\delta} = (\delta_x, \delta_y)\)
via \(|\delta_x| + |\delta_y|\) treats the two coordinates as
independently observed scalar channels.

Conversely, treating \(\vec{\delta}\) as a single vector under
continuous rotation selects the local stalk gauge
\(\|\vec{\delta}\|_2\), up to scale.

The aggregate measuring these local bounds independently across edges is
then the mixed \(\ell^1(E; \ell^2)\) measure. This work formalizes the
geometric separation between coordinate-separable scalar error
components and vector-valued Euclidean defects.

\textbf{Falsifiability.} A proposed use of this paper fails if it cannot
specify the finite carrier, vector stalk dimension \(k\ge2\), local
\(\mathrm{SO}(k)\)-invariant gauge assumption, edge-decomposable and
disjoint-additive aggregation assumptions, positive edge weights
\(c_e\), and the normalization or edge-orbit condition used to pass from
weighted to unweighted \(\ell^1(E;\ell^2)\). The duality claim fails if
a finite counterexample satisfies the primal quotient problem but
violates the stated divergence-free Euclidean-ball dual formula.

\subsubsection{1.1 Related Work and
Positioning}\label{related-work-and-positioning}

This paper is closest to five bodies of work:

\begin{enumerate}
\def\labelenumi{\arabic{enumi}.}
\item
  \textbf{Invariant norm and convex-body classification.} The
  \(\mathrm{SO}(k)\) result is a special case of norm classification
  under group actions: transitivity on the Euclidean unit sphere forces
  an invariant norm to be Euclidean up to scale (a classical fact of
  convex geometry / normed-space theory --- a finite-dimensional norm
  whose linear isometry group acts transitively on the unit sphere is
  induced by an inner product). This paper differs by applying that
  local classification to finite obstruction cochains.
\item
  \textbf{Banach-space direct sums and mixed norms.} The resulting
  diagnostic is a mixed norm: \(\ell^1\) across edges and \(\ell^2\)
  inside each vector stalk. Mixed \(\ell^p(\ell^q)\) spaces are standard
  in functional analysis (e.g., Benedek--Panzone-type mixed-norm
  spaces). This paper differs by deriving the inner/outer split from
  separate observer assumptions rather than choosing the mixed norm as a
  modeling convention.
\item
  \textbf{Finite cochain and quotient-norm obstruction theory.} The
  primal obstruction is a quotient norm modulo exact repairs. This paper
  differs by specializing the cochain geometry to
  \(\mathrm{SO}(k)\)-invariant vector stalks and deriving the
  corresponding dual witnesses.
\item
  \textbf{Convex duality / Fenchel duality and network flow.} The mixed
  duality theorem is a finite-dimensional convex-duality result: the
  dual of \(c_e\|\cdot\|_2\) is the Euclidean ball
  \(\|\varphi_e\|_2\le c_e\) by Fenchel--Rockafellar convex duality. The
  divergence-free constraint \(d_0^*\varphi=0\) makes the dual a
  vector-valued network-flow / min-cost-circulation problem. This paper
  differs by interpreting those dual variables as divergence-free vector
  circulations.
\item
  \textbf{Vector-valued sensor, strain, and displacement diagnostics.}
  Many applied systems are \emph{modelled} with Euclidean vector defects
  rather than coordinate-separable scalar errors. This paper differs by
  giving the finite observer assumptions under which that Euclidean
  local geometry is licensed. That such systems \emph{are} Euclidean is
  a modelling declaration, not a finding of this paper; the applied
  vocabulary in §1 and §5 names interpretation candidates, not realized
  physics (register row \textbf{I8}).
\end{enumerate}

These are positioning anchors. The cited external results are used only
for background context unless explicitly invoked in a theorem.

\subsubsection{1.2 Evidence and Computational
Scope}\label{evidence-and-computational-scope}

The companion claim ledger distinguishes local proofs, imported results,
and partial mechanizations. The Lean radial-collapse theorem uses a
declared ambient norm and assumes sphere constancy; the explicit planar
circle-transitivity result is a separate algebraic theorem. The finite
cycle repair kernel uses an anchored presentation, and the full
graph-incidence identification is supplied by the prose argument in §5.
Strong duality and attainment in §4 use standard finite-dimensional
convex analysis, beyond the Lean weak-duality kernel.

The executable graph engine implements nonempty finite carriers with
identity transports, common vector dimension, and scalar signed face
incidence. Its analytical cycle formula is exact as mathematics, but its
numerical evaluation and residuals use floating-point arithmetic. A
numerical \texttt{is\_optimal} flag means that the reported gap,
divergence, and ball violation pass the declared tolerance; it is not an
exact-arithmetic theorem certificate. Rotation and sphere audits test
finitely many samples and cannot establish unbounded invariance. Fresh
build and notebook replay evidence is recorded separately from retained
historical outputs in \texttt{docs/publication\_review.md}.

\subsection{2. General M-Compatible Gauge
Classification}\label{general-m-compatible-gauge-classification}

Evaluating obstruction geometries over finite carriers requires mapping
local measurement constraints onto an admissible transformation monoid.

\subsubsection{2.1 Definition}\label{definition}

Let \(V\) be a finite-dimensional normed vector space and let
\[M \subset \text{End}(V)\] be an admissible transformation monoid
representing local observer symmetries.

A faithful gauge \(j: V \to \mathbb{R}_{\geq 0}\) is
\textbf{M-compatible} if:

\begin{enumerate}
\def\labelenumi{\arabic{enumi}.}
\tightlist
\item
  \(j(x) = 0 \iff x = 0\)
\item
  \(j(\lambda x) = |\lambda| j(x)\)
\item
  \(j(x + y) \leq j(x) + j(y)\)
\item
  \(j(Tx) \leq \|T\|_{\|\cdot\|_V\to\|\cdot\|_V} j(x)\) for all
  \(T \in M\), where \(\|\cdot\|_V\) is the declared ambient norm.
\end{enumerate}

Equivalently, \(j\) is a norm whose unit ball is compatible with the
action of \(M\).

\begin{quote}
\textbf{{[}IMPORTED EVIDENCE FROM THE M-COMPATIBLE GAUGE PAPER; NOT
LOCALLY PROVED{]} Theorem 2.1 --- General M-Compatible Gauge
Classification}

Let \(V\) be finite-dimensional and \(M \subset \text{End}(V)\) an
admissible transformation monoid. There is a one-to-one correspondence
between \(M\)-compatible faithful gauges \(j\) and closed, convex,
balanced, absorbing, \textbf{bounded} unit bodies \(B \subset V\)
satisfying \[T(B) \subset \|T\| B \quad \text{for all } T \in M.\] The
correspondence is given by the Minkowski functional
\(j_B(x) = \inf\{\lambda > 0 : x \in \lambda B\}\).

Consequently, classifying M-compatible obstruction theories is
equivalent to classifying M-contractive symmetric convex bodies.
Increasing \(M\) shrinks the admissible body class. The maximal case
\(M = \text{End}_b(V)\) collapses the body class to scalar multiples of
the native Banach unit ball, recovering a unique local norm (up to
scale); \(\ell^1\) arises only after imposing replica-extensive
aggregation across edges.

This statement is reproduced solely as an imported dependency. This
paper supplies neither a local proof nor a Lean mechanization of Theorem
2.1; any use below is conditional on the cited companion source and
inherits its claim status.

Boundedness is essential and matches the companion source's
finite-dimensional statement. For example, \(B=V\) is closed, convex,
balanced, and absorbing, but its Minkowski functional is identically
zero and is not faithful. In finite dimension a closed balanced convex
absorbing body is a neighborhood of zero; boundedness excludes nonzero
directions with zero gauge.
\end{quote}

\subsubsection{2.2 Corollaries}\label{corollaries}

\begin{itemize}
\tightlist
\item
  \textbf{{[}IMPORTED FROM THE M-COMPATIBLE GAUGE PAPER / PROVEN{]}
  Corollary 1 (Order Reversal):} If \(M_1 \subset M_2\), then
  \(\text{Gauge}(V, M_2) \subset \text{Gauge}(V, M_1)\). More symmetry
  weakly reduces the admissible gauge class; strict reduction requires a
  separating example.
\item
  \textbf{{[}IMPORTED FROM THE M-COMPATIBLE GAUGE PAPER / PROVEN,
  CONDITIONAL + Hahn--Banach / native norm{]} Corollary 2 (Maximal
  Symmetry):} If \(M = \text{End}_b(V)\), then
  \(\text{Gauge}(V, M) = \{ c\|\cdot\| : c > 0 \}\). This recovers the
  maximal-local-symmetry endpoint used by the companion
  gauge-classification paper.
\item
  \textbf{{[}PROVEN, CONDITIONAL + isometric group action{]} Corollary 3
  (Group Symmetry \& Native Norm Rigidity):} If \(M=G\) is a group of
  linear isometries on a normed space \((V, \|\cdot\|_V)\), then
  \(j(gx) = j(x)\) for all \(g \in G\), and the \(M\)-compatible gauges
  are exactly the \(G\)-invariant norms. Equivalently, the unit body
  \(B_j\) is a \(G\)-invariant symmetric convex body. If \(G\) acts
  transitively on the native unit sphere \(\{x \in V : \|x\|_V = 1\}\),
  then positive homogeneity and faithfulness determine the local gauge
  up to scale as \(j(x)=c\|x\|_V\) for some \(c > 0\). In particular,
  when \(V = \mathbb{R}^k\) is equipped with the standard Euclidean
  inner product and \(G = \mathrm{SO}(k)\), this specializes to
  \(j(x) = c \|x\|_2\).
\item
  \textbf{{[}IMPORTED FROM THE M-COMPATIBLE GAUGE PAPER / STANDARD /
  PROVEN, CONDITIONAL + compact group + Haar measure{]} Corollary 4
  (Compact Group Construction):} If \(G\) is compact, acts continuously
  and linearly, and \(q\) is any norm, then
  \(j_G(x) := \int_G q(gx) dg\) is a \(G\)-invariant norm. Thus, compact
  symmetry classes generally admit vast continuous families of gauges
  unless the group acts transitively enough to collapse them.
\item
  \textbf{{[}STANDARD / SUPPORT LEMMA{]} Support Lemma S8 (Orbit Convex
  Hull Description):} For a compact group \(G\) acting continuously and
  linearly on finite-dimensional \(V\) and a bounded seed set \(S\)
  spanning \(V\), the closed balanced convex hull
  \(B=\overline{\operatorname{aconv}}(\bigcup_{g\in G}gS)\) is bounded,
  absorbing, and \(G\)-invariant. Its Minkowski functional is a faithful
  \(G\)-invariant norm. If \(S\) is already balanced, the ordinary
  closed convex hull suffices.
\end{itemize}

\emph{Warning:} There is generally \textbf{no finite list} of
\(M\)-compatible gauges. The classification maps
\(M \to \text{family of invariant/contractive convex bodies}\). Only
restrictive maximal monoids collapse the hierarchy to a single local
gauge class.

\subsection{\texorpdfstring{3. Formal Framework: Directed Graph Complex
and \(\mathrm{SO}(k)\)
Geometry}{3. Formal Framework: Directed Graph Complex and \textbackslash mathrm\{SO\}(k) Geometry}}\label{formal-framework-directed-graph-complex-and-mathrmsok-geometry}

Let \(G = (V, E, s, t)\) be a finite directed graph with vertex set
\(V\), an indexed edge set \(E\), and source and target maps
\(s,t:E\to V\). Indexed edges allow parallel comparisons with the same
endpoints. We write \(e=(u\to v)\) when \(s(e)=u\) and \(t(e)=v\).

\subsubsection{\texorpdfstring{3.1 \textbf{{[}DEFINITION{]}} The
Vector-Valued Cochain
Spaces}{3.1 {[}DEFINITION{]} The Vector-Valued Cochain Spaces}}\label{definition-the-vector-valued-cochain-spaces}

Each vertex \(v \in V\) carries a vector assignment
\(a(v) \in \mathbb{R}^k\). For each directed edge
\(e = (u \to v) \in E\), declare transport
\(r_{v,u}:\mathbb R^k_u\to\mathbb R^k_v\) in
\(\mathrm{GL}(k,\mathbb R)\). The target-stalk comparison defect is
\[\delta_e(a) = a(v) - r_{v,u}a(u)\] where
\(\delta_e \in \mathbb{R}^k\).

In the standard flat reference setting \(r_{v,u} = I\), the coboundary
operator \(d_0: C^0(G; \mathbb{R}^k) \to C^1(G; \mathbb{R}^k)\) acts on
vertex fields \(\alpha \in C^0(G; \mathbb{R}^k)\) by:
\[(d_0 \alpha)_e = \alpha(v) - \alpha(u).\]

The graph \(G = (V, E)\), the edge transports, and the field assignments
are declared combinatorial data. They provide an explicit combinatorial
index for exhibited edge comparisons rather than presupposing a
continuous ambient embedding, an antecedent physical space, or an
unindexed observer. Common over-interpretations are explicitly scoped
out:

\begin{enumerate}
\def\labelenumi{\arabic{enumi}.}
\item
  \(G\) is a declared combinatorial complex indexing comparisons, not an
  unindexed spatial continuum.
\item
  The assignment \(a\) parameterizes local vertex values; \(\delta_e\)
  is an exhibited cochain defect.
\item
  \(\mathrm{SO}(k)\) is a declared local gauge group acting on stalk
  fibers, not an external coordinate rotation of physical space.
\end{enumerate}

\subsubsection{3.2 Local Gauge Theorem}\label{local-gauge-theorem}

\textbf{{[}PROVEN, CONDITIONAL + \(k\ge2\), faithful absolute
homogeneous convex gauge{]} Theorem 3.1 --- \(\mathrm{SO}(k)\)
Invariance Forces the Euclidean Local Gauge.} Let \(k\ge2\) and let
\(j:\mathbb R^k\to\mathbb R_{\ge0}\) be a faithful absolute homogeneous
convex gauge satisfying \[
j(Rx)=j(x)
\qquad
\forall R\in\mathrm{SO}(k),\ x\in\mathbb R^k.
\] Then there is a constant \(c>0\) such that \[
j(x)=c\|x\|_2
\qquad
\forall x\in\mathbb R^k.
\]

The \(k=1\) case is the scalar one-dimensional case: every faithful
homogeneous norm is proportional to absolute value. The SO(k)
transitivity proof is only needed for \(k \ge 2\).

\emph{Proof.} The group \(\mathrm{SO}(k)\) acts transitively on every
Euclidean sphere \[
S_r^{k-1}=\{x\in\mathbb R^k:\|x\|_2=r\}
\] for \(k\ge2\). Therefore \(j\) is constant on each sphere. Define \[
\phi(r):=j(x)
\quad\text{for any }x\text{ with }\|x\|_2=r.
\] This is well-defined by \(\mathrm{SO}(k)\) invariance. By positive
homogeneity of the gauge, \[
\phi(r)=j(ru)=rj(u)
\] for any Euclidean unit vector \(u\). Since \(j\) is constant on the
unit sphere, \(j(u)=c\) for all \(\|u\|_2=1\). Faithfulness gives
\(c>0\). Hence \[
j(x)=c\|x\|_2
\] for all \(x\ne0\), and the formula also holds at \(x=0\). \(\square\)

This theorem is local. It selects the Euclidean stalk gauge only after
the observer has declared full continuous rotational symmetry.
Signed-permutation symmetry alone would not select \(\ell^2\); it admits
many unconditional norms.

The radial argument itself uses only invariance, nonnegativity,
faithfulness, and positive homogeneity; convexity and subadditivity are
not needed for this step. They remain part of the norm hypotheses used
for cycle minimization and convex duality below.

\begin{quote}
\textbf{IMPORT --- the inner product is admitted together with
\(\mathrm{SO}(k)\); it is not constructed at this degree.}

\(\mathrm{SO}(k)=\{R\in GL(k,\mathbb R): R^\top R=I,\ \det R=1\}\) is
\emph{definitionally} the isometry group of the standard inner product
on \(\mathbb R^k\): the condition \(R^\top R=I\) cannot be stated
without that inner product, and the proof above quantifies over the
Euclidean spheres \(\|x\|_2=r\) --- the selected norm appears in the
proof of its own selection. \textbf{Declaring \(\mathrm{SO}(k)\) is
therefore declaring the inner product.} Theorem 3.1 is a \emph{rigidity}
theorem relative to an already-admitted structure: a gauge declared
invariant under the isometry group of an admitted inner product cannot
differ from that inner product's norm, up to scale. It is \textbf{not} a
construction of the \(\ell^2\) fibre at this paper's declared degree,
and it does not make that fibre cheaper than the structure it arrives
with. See the \emph{Import / Oracle boundary} register, rows
\textbf{I1--I3}.

The theorem's mathematical force is real and is \textbf{not} withdrawn:
among all faithful absolutely homogeneous convex gauges on
\(\mathbb R^k\), the \(\mathrm{SO}(k)\)-invariant ones are exactly the
positive multiples of \(\|\cdot\|_2\) --- a genuine, non-vacuous
restriction of the gauge class. What is priced here is the
\emph{direction of supply}, not the validity of the implication.

The paper's own mechanization exhibits this boundary rather than hiding
it. In \texttt{lean4/gs002\_l003\_02\_s\_ok\_local\_gauge.lean}, while
an explicit constructive \(\mathrm{SO}(2)\) circle-transitivity theorem
(\texttt{so2\_transitive\_on\_circle}) is formalized for
\(\mathbb{R}^2\), the general \(k\)-dimensional theorem
\texttt{so\_k\_local\_gauge\_simplified} is stated over an arbitrary
\texttt{NormedAddCommGroup\ V}: it \emph{assumes} unit-sphere constancy
rather than constructing a general Lie group \(\mathrm{SO}(k)\) action
on \(\mathbb{R}^k\), and concludes \(j=c\|\cdot\|\) for \textbf{the
declared ambient norm} --- instantiated at \((\mathbb R^2,\ell^1)\) it
returns \(j=c\|\cdot\|_1\). The Euclidean content for general \(k\) is
carried entirely by the choice of ambient norm, i.e.~by the import. This
note records what it costs.

\subsubsection{3.3 Edge Aggregation and
Normalization}\label{edge-aggregation-and-normalization}
\end{quote}

\textbf{{[}CONDITIONAL / DEFINITIONAL AGGREGATION; local rigidity
supplied by Theorem 3.1{]} Theorem 3.2 --- Weighted Mixed Norm Under
Declared Strong Edge Aggregation.} Let \(P\) be a finite directed graph
with edge defects \(\omega_e\in\mathbb R^k\), \(k\ge2\). Suppose:

\begin{enumerate}
\def\labelenumi{\arabic{enumi}.}
\tightlist
\item
  each local gauge \(j_e\) is faithful, homogeneous, convex, and
  \(\mathrm{SO}(k)\)-invariant;
\item
  the global diagnostic is declared strongly edge-decomposable, meaning
  \(J(\omega)=\sum_{e\in E}j_e(\omega_e)\);
\item
  this declared edgewise sum is additive over disjoint replicas.
\end{enumerate}

Then there are positive edge weights \(c_e>0\) such that \[
J_{2,1,c}(\omega)=
\sum_{e\in E}c_e\|\omega_e\|_2.
\] If edges are exchange-equivalent, the weights are constant on each
edge orbit. If, in addition, a singleton-edge normalization is imposed,
\[
j_e(x)=\|x\|_2
\quad\text{on every normalized edge type,}
\] then \(c_e=1\) on those edge types and \[
J_{2,1}(\omega)=\sum_{e\in E}\|\omega_e\|_2.
\]

\emph{Proof.} By Theorem 3.1, for each edge \(e\) there exists \(c_e>0\)
such that \[
j_e(x)=c_e\|x\|_2.
\] Hypothesis (2) declares that the global diagnostic is a sum of local
edge terms: \[
J(\omega)=\sum_{e\in E}j_e(\omega_e).
\] No independent derivation of separability is claimed in this theorem.
Substitution gives the weighted mixed norm: \[
J(\omega)=\sum_{e\in E}c_e\|\omega_e\|_2.
\] If the observer declares edge-exchange equivalence, then isomorphic
edge data must receive the same singleton diagnostic, so \(c_e\) is
constant on each declared edge orbit. This gives an orbit-weighted mixed
norm. If the observer further declares singleton-edge normalization
\(j_e(x)=\|x\|_2\), then each normalized edge has \(c_e=1\). Thus the
unweighted \(\ell^1(E;\ell^2)\) form is not automatic; it is exactly the
normalized single-orbit endpoint. \(\square\)

\begin{quote}
\textbf{Remark (what is, and is not, claimed in Theorem 3.2).} Theorem
3.1 pins each local factor to \(c_e\|\cdot\|_2\). The outer summation
\(J(\omega)=\sum_e j_e(\omega_e)\) is \emph{not} a consequence of
rotational symmetry and is not newly proved here; it is the explicit
strong aggregation declaration in hypothesis (2). Support lemma
\textbf{S3} is a separate finite theorem: under its stronger zero and
disjoint-support splitting hypotheses it can derive an edgewise sum.
Weaker aggregate-statistic readings remain outside both results. Theorem
3.2 is therefore a conditional substitution statement, not an
independent aggregation theorem.

Thus, \(\mathrm{SO}(k)\) selects the local \(\ell^2\) factor, while the
declared strong edgewise rule supplies the outer \(\ell^1\) sum. The
unweighted expression is licensed only under edge normalization or
equivalent single-orbit scaling data.
\end{quote}

\subsection{4. Primal Quotient Obstruction and Analytical
Duality}\label{primal-quotient-obstruction-and-analytical-duality}

In this section and §5, transports are the identity and repairs use
\((d_0\alpha)_e=\alpha(v)-\alpha(u)\). Local adjustment fields
\(\alpha\in C^0(G,\mathbb R^k)\) determine the \textbf{Quotient
Obstruction \(\Phi_{2,1,c}\)}. For nonidentity transports the analogous
repair operator is \((d_{0,r}\alpha)_e=\alpha(v)-r_{v,u}\alpha(u)\), and
its adjoint replaces the ordinary divergence in the dual; the simple
untransported cycle-sum formula is not asserted in that setting.

\textbf{{[}DEFINITION{]} Definition 4.1 (Primal Form):} For positive
edge weights \(c_e>0\), define \[
J_{2,1,c}(\eta):=\sum_e c_e\|\eta_e\|_2
\] and \[
\Phi_{2,1,c}([\omega])
:=
\inf_{\alpha\in C^0(G,\mathbb R^k)}
\sum_e c_e\|\omega_e-(d_0\alpha)_e\|_2.
\] The normalized case is \(c_e=1\) for all \(e\), written
\(\Phi_{2,1}\).

The quotient here is \(C^1/\operatorname{im}d_0\). When faces are
present, it is not the cohomology space
\(H^1=\ker d_1/\operatorname{im}d_0\): nonclosed cochains are allowed,
and \(\Gamma\) records their closure defect. On the cycle graph without
2-cells, \(d_1=0\) and this quotient agrees with graph \(H^1\).

\begin{quote}
\textbf{{[}STANDARD CONVEX DUALITY / PROVEN FINITE APPLICATION{]}
Theorem 4.2 --- Mixed Vector Obstruction Duality} Let \(G\) be a finite
directed graph with identity transports, edge defects in
\(\mathbb{R}^k\), and positive edge weights \(c_e>0\). The primal
quotient obstruction satisfies: \[
\Phi_{2,1,c}([\omega])
=
\max_\varphi \sum_e \langle \varphi_e, \omega_e \rangle
\] Subject to explicit structural bounds:

\begin{enumerate}
\def\labelenumi{\arabic{enumi}.}
\tightlist
\item
  Gauge Invariance / Divergence-Free Constraint: \(d_0^* \varphi = 0\)
\item
  Weighted Euclidean Witness Bounds: \(\|\varphi_e\|_2 \le c_e\) for all
  edges \(e \in E\).
\end{enumerate}

\emph{Proof.} Write the primal problem as \[
\inf_{\alpha}\sum_e c_e\|\omega_e-(d_0\alpha)_e\|_2.
\] The convex conjugate of the function \[
x\mapsto c_e\|x\|_2
\] is the indicator function of the closed Euclidean ball \[
\{\varphi_e:\|\varphi_e\|_2\le c_e\}.
\] Using \[
c_e\|r_e\|_2=
\max_{\|\varphi_e\|_2\le c_e}\langle\varphi_e,r_e\rangle,
\] the primal becomes \[
\inf_\alpha
\max_{\|\varphi_e\|_2\le c_e}
\sum_e
\left\langle\varphi_e,\omega_e-(d_0\alpha)_e\right\rangle.
\] Because all spaces are finite-dimensional, the objective is proper,
convex, and continuous, and the primal is feasible. Fenchel--Rockafellar
duality permits exchanging infimum and supremum. The
\(\alpha\)-dependent term is \[
-\left\langle d_0^*\varphi,\alpha\right\rangle.
\] The infimum over unconstrained \(\alpha\) is finite exactly when \[
d_0^*\varphi=0.
\] Otherwise the infimum is \(-\infty\). Therefore admissible dual
witnesses are precisely divergence-free vector circulations with
weighted Euclidean bounds. Thus \[
\Phi_{2,1,c}([\omega])=
\max_{\substack{d_0^*\varphi=0\\ \|\varphi_e\|_2\le c_e}}
\sum_e\langle\varphi_e,\omega_e\rangle.
\] \(\square\)
\end{quote}

Both extrema are attained. The repair image is a closed
finite-dimensional subspace. Positive weights make \(J_{2,1,c}\) a norm,
so its sublevel sets on the affine residual class are compact and
contain a minimizer; a repair field exists for that residual, although
it need not be unique. The dual feasible set is a closed bounded
intersection of edge balls with \(\ker d_0^*\), contains zero, and is
compact. Its linear objective therefore attains a maximum. These
attainment arguments supplement the continuity qualification used for
strong duality above.

In the normalized case \(c_e=1\), the dual constraints reduce to \[
\|\varphi_e\|_2\le1.
\] Thus unweighted Euclidean-unit-bounded circulations are exactly the
edge-normalized endpoint, not the general weighted case.

\subsection{5. Finite Triangle-Anomaly
Separation}\label{finite-triangle-anomaly-separation}

To isolate the observer-level distinction, we identify a concrete finite
cycle defect.

Consider a minimal closed triangle network \(V=\{0,1,2\}\) with directed
edges \(E=\{(0,1),(1,2),(2,0)\}\). Let the defect signal over the edges
be:

\begin{itemize}
\tightlist
\item
  \(\omega_{01} = (1,0)\)
\item
  \(\omega_{12} = (0,1)\)
\item
  \(\omega_{20} = (0,0)\)
\end{itemize}

The total cycle circulation evaluates to:
\[c = \omega_{01} + \omega_{12} + \omega_{20} = (1,1)\]

\textbf{{[}STANDARD --- graph cycle-space / \(H^1\) of a single cycle;
finite exactness kernel Lean-mechanized{]}} For a single directed cycle
with identity restrictions, the quotient class is determined by the
total circulation \[
S(\omega)=\sum_{e\in C}\omega_e.
\] Indeed, exact repairs \(d_0\alpha\) have zero cycle sum, and any two
cochains with the same cycle sum differ by an exact repair (the cochains
of fixed circulation form a single coset of \(\operatorname{im} d_0\) on
a one-cycle graph).

Explicitly, label the cycle edges \(e_i=(v_i\to v_{i+1})\), with vertex
indices modulo \(n\). For \(q=\omega-\eta\) with \(\sum_i q_{e_i}=0\),
set \(\alpha(v_0)=0\) and \(\alpha(v_i)=\sum_{r=0}^{i-1}q_{e_r}\) for
\(1\le i<n\). Then \((d_0\alpha)_{e_i}=q_{e_i}\) on every nonclosing
edge, and the closing equation follows from the zero sum. If stored edge
directions are reversed, use the signed traversal circulation
\(\sum_i\varepsilon_i\omega_{e_i}\) with \(\varepsilon_i\in\{\pm1\}\);
reorienting the cochain reduces to the same construction.

The finite algebraic core of this bridge is mechanized in Lean as
\texttt{same\_circulation\_iff\_difference\_anchored\_cycle} and
\texttt{same\_circulation\_iff\_anchored\_cycle\_coset}: after choosing
an anchored cycle presentation with one closing edge, two cochains have
the same circulation exactly when their pointwise difference is an
anchored zero-circulation cycle repair, equivalently when they lie in
the same anchored repair coset. The remaining prose-standard step is
only the identification of that anchored presentation with the usual
graph coboundary notation for a concrete cycle.

\textbf{{[}PROVEN, CONDITIONAL + single cycle + normalized equal edge
weights + norm gauge{]} Lemma 5.1 (Cycle quotient norm).} Let \(C\) be a
finite directed cycle with edge-local norm \(j\) and normalized edge
weights. For any vector-valued cochain \(\omega\), \[
\Phi_j([\omega])=j(S(\omega)).
\]

\emph{Proof.} Let \(c=S(\omega)\). Every representative \(\eta\) of the
same quotient class satisfies \(S(\eta)=c\). By the triangle inequality,
\[
\sum_{e\in C}j(\eta_e)\ge j\left(\sum_{e\in C}\eta_e\right)=j(c).
\] Conversely, place all circulation on a single edge: choose
\(\eta_{e_0}=c\) and \(\eta_e=0\) for \(e\ne e_0\). This representative
has cycle sum \(c\) and cost \(j(c)\). Therefore the quotient infimum
equals \(j(c)\). \(\square\)

The normalized normed-additive kernel of this argument is mechanized in
Lean as \texttt{normalized\_cycle\_quotient}.

\textbf{{[}PROVEN, CONDITIONAL + single cycle + positive weights{]}
Weighted refinement.} If the edge weights are \(w_e>0\), the same
argument gives \[
\Phi_{j,w}([\omega])=
\left(\min_{e\in C}w_e\right)j(S(\omega)).
\] The lower bound follows from
\(\sum_e w_e j(\eta_e)\ge w_{\min}\sum_e j(\eta_e)\ge w_{\min}j(S(\eta))\),
and the upper bound is achieved by placing all circulation on an edge of
minimal weight. Thus clean numerical comparison requires equal
normalized edge weights. The weighted finite-cycle refinement is
mechanized in Lean as \texttt{weighted\_cycle\_quotient}.

\textbf{{[}PROVEN, NUMERICAL / DIRECT COMPUTATION{]} Theorem 5.2 (Finite
triangle-anomaly separation).} On the normalized triangle above, \[
\Phi_{1,1}([\omega])=2,
\qquad
\Phi_{2,1}([\omega])=\sqrt2.
\] Hence coordinate-separable scalar observers and
\(\mathrm{SO}(2)\)-invariant vector observers assign different
obstruction magnitudes to the same finite cycle class.

\begin{quote}
\textbf{Algebraic Extension Note on \(\sqrt{2}\).}

The two magnitudes are on distinct scalar footings. \(\Phi_{1,1} = 2\)
is a rational scalar obtained by summing integer-valued cochain
components. \(\Phi_{2,1} = \sqrt{2}\) is irrational and requires at
minimum the quadratic field extension \(\mathbb{Q}(\sqrt{2})\) over the
rationals.

In the formalization in
\texttt{lean4/gs002\_l003\_01\_triangle\_separation.lean}, this boundary
is verified rationally by proving the squared inequality \(2^2 \ne 2\)
(that is, \(4 \ne 2\)) in \(\mathbb{Q}\), certifying that the coordinate
\(\ell^1\) cost and \(\mathrm{SO}(2)\) Euclidean cost are separated
without requiring the non-constructive continuum square root.

Minimal replay:
\end{quote}

\begin{Shaded}
\begin{Highlighting}[]
\ImportTok{import}\NormalTok{ math}

\NormalTok{omega01 }\OperatorTok{=}\NormalTok{ (}\DecValTok{1}\NormalTok{, }\DecValTok{0}\NormalTok{)}
\NormalTok{omega12 }\OperatorTok{=}\NormalTok{ (}\DecValTok{0}\NormalTok{, }\DecValTok{1}\NormalTok{)}
\NormalTok{omega20 }\OperatorTok{=}\NormalTok{ (}\DecValTok{0}\NormalTok{, }\DecValTok{0}\NormalTok{)}

\NormalTok{c }\OperatorTok{=} \BuiltInTok{tuple}\NormalTok{(a }\OperatorTok{+}\NormalTok{ b }\OperatorTok{+}\NormalTok{ d }\ControlFlowTok{for}\NormalTok{ a, b, d }\KeywordTok{in} \BuiltInTok{zip}\NormalTok{(omega01, omega12, omega20))}

\NormalTok{phi\_11 }\OperatorTok{=} \BuiltInTok{abs}\NormalTok{(c[}\DecValTok{0}\NormalTok{]) }\OperatorTok{+} \BuiltInTok{abs}\NormalTok{(c[}\DecValTok{1}\NormalTok{])}
\NormalTok{phi\_21 }\OperatorTok{=}\NormalTok{ math.sqrt(c[}\DecValTok{0}\NormalTok{]}\OperatorTok{**}\DecValTok{2} \OperatorTok{+}\NormalTok{ c[}\DecValTok{1}\NormalTok{]}\OperatorTok{**}\DecValTok{2}\NormalTok{)}

\BuiltInTok{print}\NormalTok{(c)}
\BuiltInTok{print}\NormalTok{(phi\_11)}
\BuiltInTok{print}\NormalTok{(phi\_21)}
\end{Highlighting}
\end{Shaded}

Expected output:

\begin{Shaded}
\begin{Highlighting}[]
\NormalTok{(1, 1)}
\NormalTok{2}
\NormalTok{1.4142135623730951}
\end{Highlighting}
\end{Shaded}

This replay checks only the normalized triangle quotient values. It does
not certify the general duality theorem or the replica-extensive
aggregation theorem.

\emph{Proof.} The cycle sum is \[
c=(1,0)+(0,1)+(0,0)=(1,1).
\] For the coordinate-separable scalar observer, the local gauge is \[
j_{1,1}(x,y)=|x|+|y|.
\] By Lemma 5.1, \[
\Phi_{1,1}([\omega])=j_{1,1}(c)=|1|+|1|=2.
\] For the \(\mathrm{SO}(2)\)-invariant vector observer, the local gauge
is \[
j_{2}(x,y)=\|(x,y)\|_2.
\] Again by Lemma 5.1, \[
\Phi_{2,1}([\omega])=j_2(c)=\sqrt{1^2+1^2}=\sqrt2.
\] Since \(2\ne\sqrt2\), the finite triangle separates the observer
classes. \(\square\)

\subsubsection{\texorpdfstring{Prediction A --- Scalar-Coordinate
\(\ell^1\)
(\(J_{1,1}\))}{Prediction A --- Scalar-Coordinate \textbackslash ell\^{}1 (J\_\{1,1\})}}\label{prediction-a-scalar-coordinate-ell1-j_11}

If the topological observer separates variables (e.g.~tracking
independent coordinate channels):
\(J_{1,1}(\omega) = \sum_e (|\omega_{e,x}| + |\omega_{e,y}|)\) yields a
quotient obstruction of: \[\Phi_{1,1} = \|c\|_1 = |1| + |1| = 2\]

\subsubsection{\texorpdfstring{Prediction B --- Vector Mixed Norm
(\(J_{2,1}\))}{Prediction B --- Vector Mixed Norm (J\_\{2,1\})}}\label{prediction-b-vector-mixed-norm-j_21}

If the topological observer operates under continuous \(\mathrm{SO}(k)\)
rotational invariance: \(J_{2,1}(\omega) = \sum_e \|\omega_e\|_2\)
yields a quotient obstruction of:
\[\Phi_{2,1} = \|c\|_2 = \sqrt{1^2 + 1^2} = \sqrt{2} \approx 1.414\]

\subsubsection{Finite Numerical
Consequence}\label{finite-numerical-consequence}

\textbf{{[}PROVEN, CONDITIONAL on normalization{]}} The two observers
assign different quotient norms to the same cycle class.

An observer declaring separable coordinates (\(J_{1,1}\)) --- the
\emph{component-counting} reading --- assigns the class magnitude
\textbf{2}. An observer declaring \(\mathrm{SO}(k)\)-invariant vector
stalks (\(J_{2,1}\)) --- the \emph{Euclidean-distance} reading ---
assigns the same class magnitude \textbf{\(\sqrt{2}\)}.

Both numbers are computed from declared gauges on declared data. Naming
the first a ``component-counting anomaly detector'' and the second a
``structural strain sensor'' describes \emph{candidate interpretations}
of the two observer classes; neither is an instrument reading, neither
measures a physical system, and no claim is made that any physical
detector realizes either gauge (register row \textbf{I8}).

This distinguishes observer classes at the level of local gauge choice;
it does not by itself test replica-extensive aggregation or dual
structure.

The weighted refinement is important. If the scalar-coordinate observer
uses edge weights \(a_e\) and the \(\mathrm{SO}(2)\) observer uses edge
weights \(b_e\), then the triangle values become \[
\Phi_{1,1,a}=2\min_e a_e,
\qquad
\Phi_{2,1,b}=\sqrt2\min_e b_e.
\] Thus the canonical \(2\) versus \(\sqrt2\) separation is a normalized
comparison. Without normalization, the theorem shows observer-class
dependence together with explicit scale dependence, not an absolute
numerical invariant.

\subsection{6. Closure Invariant Bound
Extension}\label{closure-invariant-bound-extension}

Under topological evaluation, the 2-cell failure metric inherits the
declared mixed geometry. Let \[
C^0\xrightarrow{d_0}C^1\xrightarrow{d_1}C^2,
\qquad
d_1d_0=0
\] be a finite obstruction complex with
\(C^1=\bigoplus_{e\in E}\mathbb R^k\) and
\(C^2=\bigoplus_{f\in F}\mathbb R^{m_f}\). Equip \(C^1\) with \[
\|\eta\|_{1,2,c}:=\sum_{e\in E}c_e\|\eta_e\|_2
\] and equip \(C^2\) with a declared mixed closure norm \[
\|\zeta\|_{1,2,u}:=\sum_{f\in F}u_f\|\zeta_f\|_2,
\qquad
c_e,u_f>0.
\]

Define \[
\Phi_{2,1,c}([h])
:=
\inf_\alpha \|h-d_0\alpha\|_{1,2,c}
\] and \[
\Gamma_{2,1,u}([h])
:=
\|d_1h\|_{1,2,u}.
\] The closure diagnostic is well-defined on quotient classes because
\(d_1d_0=0\).

\textbf{{[}PROVEN, CONDITIONAL + finite complex + \(d_1d_0=0\) +
declared mixed norms{]} Theorem 6.1 --- Mixed Closure Bound.} With the
definitions above, \[
\Gamma_{2,1,u}([h])
\le
\|d_1\|_{(1,2,c)\to(1,2,u)}
\,
\Phi_{2,1,c}([h]).
\]

\emph{Proof.} Since \(d_1d_0=0\), for every repair field \(\alpha\), \[
d_1(h-d_0\alpha)=d_1h.
\] By the operator norm definition, \[
\|d_1h\|_{1,2,u}=
\|d_1(h-d_0\alpha)\|_{1,2,u}
\le
\|d_1\|_{(1,2,c)\to(1,2,u)}
\,
\|h-d_0\alpha\|_{1,2,c}.
\] Taking the infimum over all \(\alpha\) gives \[
\Gamma_{2,1,u}([h])
\le
\|d_1\|_{(1,2,c)\to(1,2,u)}
\Phi_{2,1,c}([h]).
\] \(\square\)

In the normalized case \(c_e=u_f=1\), this reduces to the unweighted
bound \[
\Gamma_2([h])
\le
\|d_1\|_{(2,1\to2,1)}
\Phi_{2,1}([h]).
\] This shows the mixed closure bound --- the \(\Phi\to\Gamma\) link of
the \(\Phi\Gamma R_{\mathrm{cl}}\) flow --- holds for Euclidean
rotational fields, with the operator constant carrying the declared edge
and face normalization data. The closure ratio \(R_{cl}=\Gamma/\Phi\) is
defined only when \(\Phi>0\). When \(\Phi=0\), the ratio is not used;
the bound is the \(\Phi\to\Gamma\) inequality itself.

For the implementation's common stalk dimension and scalar signed
incidence matrix \(A=d_1\), the operator constant has the explicit
formula \[
\|d_1\|_{(1,2,c)\to(1,2,u)}
=\max_{e\in E}\frac{\sum_{f\in F}u_f|A_{fe}|}{c_e}.
\] The triangle inequality bounds each face norm by
\(\sum_e|A_{fe}|\|h_e\|_2\), giving the upper bound after weighted
summation. A cochain supported on a maximizing edge attains it. This
formula concerns scalar incidence acting on equal-dimensional stalks;
the abstract theorem also allows the more general declared linear
operator and face dimensions specified above.

\subsection{7. Conclusion}\label{conclusion}

Continuous \(\mathrm{SO}(k)\) invariance makes the local gauge Euclidean
up to scale. The separately declared edgewise sum supplies the outer
\(\ell^1\) form, with weights fixed by singleton normalization or other
declared scaling data. Edge exchange symmetry alone supplies common
weights on each orbit. In the flat finite graph setting, the resulting
quotient norm has the stated dual witnesses, single-cycle formula, and
closure bound. The contribution is this explicit separation of
hypotheses and their finite obstruction consequences.

\subsubsection{7.1 Connection to General Obstruction
Theory}\label{connection-to-general-obstruction-theory}

The \(\ell^1(E; \ell^2)\) geometry derived here integrates into general
finite obstruction theory as a specific gauge selection within the
broader family of \(M\)-compatible gauges:

\textbf{Observer passivity.} In general obstruction models, observers
are contractive quotient maps satisfying
\(\Phi_{\mathcal O}(\mathcal O(q)) \le \Phi(q)\). When \(\Phi\) is
computed using the weighted or normalized \(\ell^1(E; \ell^2)\)
geometry, the observer passivity bound constrains how
\(\mathrm{SO}(k)\)-measured defects can be processed by downstream
observer maps: no observer quotient can create obstruction absent from
the rotationally gauged cochain core.

\textbf{Contextual distinction.} The testable triangle prediction (§5)
should be read as an observer-class separation: separable-coordinate and
\(\mathrm{SO}(k)\)-invariant observers assign different quotient
magnitudes to the same cycle class (\(2\) versus \(\sqrt{2}\)). It
illustrates that quotient defect magnitudes depend essentially on
declared local measurement symmetry rather than reflecting a single
observer-independent metric.

\textbf{Sheaf and network gluing.} The target-stalk defect
\(\delta_e(a)=a(v)-r_{v,u}a(u)\) represents local mismatch. Declared
\(\mathrm{SO}(k)\)-invariant local gauges and the strong edgewise sum
give the weighted or normalized \(\ell^1(E;\ell^2)\) diagnostic. Theorem
4.2 supplies ordinary divergence-free witnesses for identity transports.
With nonidentity transports, dual witnesses instead satisfy the adjoint
constraint for \(d_{0,r}\); no untwisted circulation formula or holonomy
lower bound is supplied here.

\newpage

\subsection{References}\label{references}

{[}1{]} J. H. Carroll, ``Primitive Axiomatics of Finite Obstruction
Calculus: Evaluation Preorders, Quotient Residuals, Closure Charge, and
Observer-Selected Geometry,'' preprint, 2026.
\href{https://github.com/dbccpage/the_symbol_game/tree/adc3d57001a0491945fce1ccd203f696c57c389b/omega_engine_atlas/coda_atlas/F004_papers/001_finite_obstruction_calculus/006_primitive_axiomatics}{Repository
source}.

{[}2{]} J. H. Carroll, ``M-Compatible Gauge Classification for Banach
Presheaf Coboundaries: The ℓ¹ Maximal-Symmetry Endpoint,'' preprint,
2026. \href{https://zenodo.org/records/20677155}{Zenodo record
20677155}.
\href{https://github.com/dbccpage/the_symbol_game/tree/adc3d57001a0491945fce1ccd203f696c57c389b/omega_engine_atlas/coda_atlas/F004_papers/002_projection_obstruction/007_m_compatible_gauge_classification}{Companion
source at the reviewed commit}, finite-dimensional unit-body
correspondence, §2.2, Theorem 2.2.

{[}3{]} J. H. Carroll, ``Projection Obstruction for Factorized States:
Nonlinear Dynamics, Conditional ℓ1 Diagnostics, and Finite-Time Defect
Bounds,'' preprint, 2026.
\href{https://zenodo.org/records/21752743}{Zenodo record 21752743}.

{[}4{]} J. H. Carroll, ``Coordinate-Wise Additivity and the ℓ¹ Norm on
Finite Graph Cochains,'' preprint, 2026.
\href{https://zenodo.org/records/20723753}{Zenodo record 20723753}.

{[}5{]} J. H. Carroll, ``L1-Native Cohomology of Finite Defect Fields:
Quotient Norms, Dual Circulations, Sparse Representatives, and Integral
Optima,'' preprint, 2026.
\href{https://zenodo.org/records/21783292}{Zenodo record 21783292}.

{[}6{]} R. T. Rockafellar, \emph{Convex Analysis}. Princeton University
Press, 1970. (For invariant norms, convex-body classification, and
Fenchel convex duality).

{[}7{]} A. Benedek and R. Panzone, ``The space \(L^p\), with mixed
norm,'' \emph{Duke Math. J.} 28(3): 301-324, 1961. (For mixed-norm
spaces).

{[}8{]} R. K. Ahuja, T. L. Magnanti, and J. B. Orlin, \emph{Network
Flows: Theory, Algorithms, and Applications}. Prentice Hall, 1993. (For
network-flow / circulation duality).
\end{document}
